% Integration fragment. Requires amsmath, amssymb, amsthm and hyperref.
% All labels and bibliography keys have the prefix mtp:.
\section{A polarized Bernoulli calculus for arbitrary mixed even tails}
\label{mtp:section}

The all-power continuation in \emph{Bernoulli Tails and Spectral Faces}
proves its recurrence for powers of one even-order Hurwitz tail and
explicitly leaves products containing unequal orders among three or
more factors outside its scope \cite[Section 2, after the even-order
extension]{RepoBernoulli}. The earlier mixed-moment theorem treats
two factors \cite[Section 4]{RepoMixed}. We now remove both
restrictions: the tail orders may be arbitrary positive even integers,
and the number of factors is unrestricted. The result is one finite
triangular recurrence, an anchored analytic generating function, and a
fixed finite arithmetic space for all regular centered moments.

The unweighted sum of an arbitrary product of zeta tails, and its
ordered-index interpretation in multiple zeta values, are classical
\cite{Hoffman}. Those sums are used below as explicit anchors.
The additional assertion concerns \emph{every} centered polynomial
weight, including weights requiring polynomial subtraction. The proof
keeps the cutoff finite until its boundary terms have cancelled.
In particular, this is a continuation of a specified repository question;
we do not claim that the classical tail-product anchors are new.

\subsection{Definitions and ordinary realization of the regular moments}

Let $I$ be a nonempty finite set of labels, let $d=|I|$, and assign an
even integer $p_i\ge2$ to each $i\in I$. Equal assigned integers retain
their different labels. Write
\[
 P_S=\sum_{i\in S}p_i,\qquad w_I=P_I-d,\qquad
 T_p(n)=\zeta(p,n+1)=\zeta(p)-H_n^{(p)},
 \qquad F_I(n)=\prod_{i\in I}T_{p_i}(n),
\]
and put $F_\varnothing(n)=1$, $x_n=n+1/2$. The polygamma normalization is
$T_p(n)=\psi^{(p-1)}(n+1)/(p-1)!$ for even $p$.

For $j\ge0$ define
\begin{equation}\label{mtp:eta}
 \eta_{p,0}=\frac1{p-1},\qquad
 \eta_{p,j}=\frac{(p)_{2j-1}}{(2j)!}B_{2j}(1/2)\quad(j\ge1),
 \qquad
 h_{I,j}=\sum_{\substack{j_i\ge0\\\sum_i j_i=j}}
                        \prod_{i\in I}\eta_{p_i,j_i}.
\end{equation}
The centered Hurwitz asymptotics \cite[Section 25.11]{DLMF} give
\begin{equation}\label{mtp:asymptotic}
                F_I(n)\sim\sum_{j\ge0}h_{I,j}x_n^{-w_I-2j}.
\end{equation}
Every $\eta_{p,j}$ has sign $(-1)^j$, including the positive $j=0$
coefficient. Consequently $h_{I,j}$ is nonzero with the same sign.
This is an asymptotic expansion, not a convergent Laurent series at a
fixed finite $n$.

\begin{proposition}[Regular moments as convergent sums]
\label{mtp:ordinary}
The Dirichlet series
\[
                 \mathcal F_I(s)=\sum_{n\ge0}F_I(n)x_n^{-s},
                    \qquad \Re s>1-w_I,
\]
continues meromorphically to the plane. Its poles are precisely
$1-w_I-2j$, with residue $h_{I,j}$, for $j\ge0$. For every integer
$m\ge0$ satisfying $m\equiv d\pmod2$, its regular value
$M_{I,m}=\mathcal F_I(-m)$ is the ordinary absolutely convergent sum
\begin{equation}\label{mtp:ordinary-formula}
 \boxed{M_{I,m}=\sum_{n=0}^{\infty}\left\{
       x_n^mF_I(n)
       -\sum_{0\le j\le(m-w_I)/2}h_{I,j}x_n^{m-w_I-2j}
                                     \right\}.}
\end{equation}
A sum with a negative upper limit is empty.
\end{proposition}

\begin{proof}
Subtract finitely many terms of \eqref{mtp:asymptotic} from the defining
series and add back the corresponding Hurwitz functions
$h_{I,j}\zeta(s+w_I+2j,1/2)$. Increasing the subtraction order produces
normally convergent remainders in successively larger half-planes.
The nonzero coefficients prove that all the listed poles are actual.
Since $w_I\equiv d\pmod2$, none of these poles is $-m$ on the stated
parity. Every subtracted term in \eqref{mtp:ordinary-formula} contributes
$\zeta(-2r,1/2)=0$, $r\ge0$. After the subtraction the summand is
$O(x_n^{-2})$, or decays faster when no subtraction is needed.
This proves both the value and its ordinary convergence.
\end{proof}

For a full asymptotic expansion in the integer cutoff $N$, let
$\operatorname{CT}_N$ mean the coefficient of $N^0(\log N)^0$.
For any integer $k$ put
\begin{equation}\label{mtp:C}
 C_{I,k}=\operatorname{CT}_N\left\{
     \sum_{n=1}^N n^k\bigl(F_I(n-1)-F_I(n)\bigr)+N^kF_I(N)\right\}.
\end{equation}
These coefficients exist by asymptotic subtraction: each tail has an
expansion in integral inverse powers of $N$, and finite sums of such
powers have power--logarithm expansions with a constant coefficient.
The convention for this coefficient is fixed by \eqref{mtp:C}.

The centered Faulhaber polynomial is
\[
 P_m(N)=\sum_{n=0}^{N-1}(n+1/2)^m
       =\frac{B_{m+1}(N+1/2)-B_{m+1}(1/2)}{m+1}.
\]
Finite summation by parts gives
\[
 \sum_{n=0}^{N-1}x_n^mF_I(n)
 =P_m(N)F_I(N)+\sum_{n=1}^NP_m(n)\bigl(F_I(n-1)-F_I(n)\bigr).
\]
All the removed even-power sums in \eqref{mtp:ordinary-formula} are odd
polynomials in $N$, with zero constant coefficient. Therefore
\begin{equation}\label{mtp:moment-transform}
 M_{I,m}=\sum_{k=1}^{m+1}
       \frac{\binom{m+1}{k}}{m+1}B_{m+1-k}(1/2)C_{I,k}.
\end{equation}
Only indices $k\equiv d+1\pmod2$ occur. Define initially as formal series
\begin{equation}\label{mtp:GH}
 G_I(t)=\sum_{\substack{m\ge0\\m\equiv d\ (2)}}M_{I,m}\frac{t^m}{m!},
 \qquad H_I(t)=2\sinh(t/2)G_I(t),\qquad G_\varnothing=H_\varnothing=0.
\end{equation}
The Faulhaber generating function makes \eqref{mtp:moment-transform}
equivalent to
\begin{equation}\label{mtp:HC}
 H_I(t)=\sum_{\substack{k\ge1\\k\equiv d+1\ (2)}}C_{I,k}\frac{t^k}{k!}.
\end{equation}
No interchange with a divergent infinite moment sum is used.

\subsection{The finite polarized identity and the exact boundary parity}

For a polynomial in independent variables $a_i$, write
$\mathscr D=\sum_i h_i\partial_{a_i}$. On a multilinear polynomial this
operator is nilpotent. The scalar Bernoulli identity
\[
 e^z-1=\frac z2(e^z+1)
          -\sum_{r\ge1}\frac{B_{2r}}{(2r)!}z^{2r}(e^z-1)
\]
can therefore be applied to $\prod_{i\in I}a_i$ as a finite polynomial
identity. It is equivalent to the usual generating function
$(z/2)\coth(z/2)=1+\sum_{r\ge1}B_{2r}z^{2r}/(2r)!$.

\begin{lemma}[Polarized Bernoulli difference]
\label{mtp:polarized}
If $\Delta F_J(n)=F_J(n-1)-F_J(n)$, then for every $n\ge1$
\begin{equation}\label{mtp:polarized-formula}
 \boxed{\begin{aligned}
 \Delta F_I(n)={}&\frac12\sum_{i\in I}n^{-p_i}
                         \bigl(F_{I\setminus\{i\}}(n-1)
                              +F_{I\setminus\{i\}}(n)\bigr)\\
 &-\sum_{\substack{S\subsetneq I\\|S|\ge2\text{ even}}}
                   B_{|S|}n^{-P_S}\Delta F_{I\setminus S}(n).
 \end{aligned}}
\end{equation}
Every product on the right contains fewer tail factors.
\end{lemma}

\begin{proof}
Apply the displayed scalar identity with
$a_i=T_{p_i}(n)$ and $h_i=n^{-p_i}$. Since
$T_{p_i}(n-1)=T_{p_i}(n)+n^{-p_i}$, the exponential of $\mathscr D$
implements the simultaneous shift. The identity
\[
 \frac{\mathscr D^r}{r!}\prod_{i\in I}a_i
  =\sum_{\substack{S\subseteq I\\|S|=r}}
                  \left(\prod_{i\in S}h_i\right)\prod_{i\notin S}a_i
\]
accounts for every coefficient. The term $S=I$ would be the difference
of two constants and is zero. Labels, rather than distinct numerical
orders, supply the multiplicities when some $p_i$ coincide.
\end{proof}

Let $\mathcal A_p(u)$ be the formal expansion of $T_p(N)$ at $u=N^{-1}$:
\begin{equation}\label{mtp:unshifted}
 \mathcal A_p(u)=\frac{u^{p-1}}{p-1}-\frac{u^p}{2}
       +\sum_{r\ge1}\frac{B_{2r}(p)_{2r-1}}{(2r)!}u^{p+2r-1}.
\end{equation}
Put $b_I(k)=[u^k]\prod_{i\in I}\mathcal A_{p_i}(u)$, with
$b_\varnothing(k)=\delta_{k,0}$, and define the odd formal series
$E_i(u)=\mathcal A_{p_i}(u)+u^{p_i}/2$.

\begin{lemma}[Boundary cancellation for unequal increments]
\label{mtp:boundary}
For $k\equiv d+1\pmod2$,
\begin{equation}\label{mtp:boundary-coefficient}
 b_I(k)+\frac12\sum_{i\in I}b_{I\setminus\{i\}}(k-p_i)
 +\sum_{\substack{S\subseteq I\\|S|\ge2\text{ even}}}
                    B_{|S|}b_{I\setminus S}(k-P_S)=0.
\end{equation}
More precisely, the entire generating expression on the left is
\begin{equation}\label{mtp:boundary-polynomial}
 \boxed{\sum_{\substack{S\subseteq I\\|S|\text{ even}}}
             B_{|S|}(1/2)u^{P_S}\prod_{i\notin S}E_i(u).}
\end{equation}
It has parity $d$.
\end{lemma}

\begin{proof}
Let $B_r^+=B_r(1)$, so $B_1^+=1/2$, $B_r^+=B_r$ for $r\ne1$,
and odd indices above one vanish. Before taking a coefficient the
left side of \eqref{mtp:boundary-coefficient} is
\[
                 \sum_{S\subseteq I}B_{|S|}^+u^{P_S}
                              \prod_{i\notin S}\mathcal A_{p_i}(u).
\]
Substitute $\mathcal A_{p_i}=E_i-u^{p_i}/2$. For a fixed subset $R$,
the coefficient of $u^{P_R}\prod_{i\notin R}E_i$ is
\[
 \sum_{S\subseteq R}B_{|S|}(1)(-1/2)^{|R|-|S|}
      =B_{|R|}(1/2),
\]
by the Bernoulli translation identity \cite[24.4.12]{DLMF}.
Odd $|R|$ vanish. Each $u^{P_R}$ is even and the remaining product
has parity $d-|R|\equiv d$. This proves the identity and cancellation.
\end{proof}

This is the point where equal tail orders are unnecessary. Each tail
has its own even increment $u^{p_i}$; Bernoulli translation cancels the
entire mixed boundary polynomial at once. If $d$ is even, the full
subset in \eqref{mtp:boundary-coefficient} has a possible term only at
$k=P_I$, an even index outside the required odd parity. Thus its
inclusion is harmless, although the corresponding difference in
\eqref{mtp:polarized-formula} is omitted.

\subsection{Every anchor is a specified convergent multiple zeta value}

Our convention is
\[
 \zeta(s_1,\ldots,s_\ell)
   =\sum_{n_1>\cdots>n_\ell\ge1}n_1^{-s_1}\cdots n_\ell^{-s_\ell}.
\]
An ordered set partition of $J$ is a sequence
$\pi=(B_1,\ldots,B_\ell)$ of disjoint nonempty blocks whose union is
$J$. For $r\ge1$ and nonempty $J$ define
\begin{align}
 \mathcal T_{J,r}
   &=\sum_{\pi\models J}\zeta(P_{B_1},\ldots,P_{B_\ell},r),
       &\mathcal T_{\varnothing,r}&=\zeta(r)\quad(r\ge2),
       \label{mtp:Tanchor}\\
 \mathcal S_{J,r}
   &=2\mathcal T_{J,r}
       +\sum_{\varnothing\ne U\subseteq J}\mathcal T_{J\setminus U,r+P_U}.
       \label{mtp:Sanchor}
\end{align}
The only use of $J=\varnothing$ below has $r$ even and at least two.
For nonempty $J$ and $h\le0$ put
\begin{equation}\label{mtp:Zanchor}
 \mathcal Z_{J,h}=\sum_{\pi\models J}
       \zeta(P_{B_1},\ldots,P_{B_{\ell-1}},P_{B_\ell}-h),
 \qquad Z_J=\prod_{i\in J}\zeta(p_i),\quad Z_\varnothing=1.
\end{equation}
All these MZVs are convergent and have positive integer entries with
first entry at least two. In particular no isolated $\zeta(1)$ is
introduced.

Index ordering, or direct partition according to the indices equal to
$n$, proves the ordinary identities
\begin{align}
 \mathcal T_{J,r}&=\sum_{n\ge1}n^{-r}F_J(n),\notag\\
 \mathcal S_{J,r}&=\sum_{n\ge1}n^{-r}\bigl(F_J(n-1)+F_J(n)\bigr),
       \label{mtp:anchor-series}\\
 \mathcal Z_{J,h}&=\sum_{n\ge1}n^h\Delta F_J(n)=C_{J,h}\quad(h\le0).
       \notag
\end{align}
The last equality uses $N^hF_J(N)\to0$. These are instances of the
classical tail-index method \cite{Hoffman}; their explicit form
makes every constant in the new recurrence independently computable.
Their respective weights are $P_J+r$, $P_J+r$, and $P_J-h$.

\subsection{The all-multiorder recurrence}

For $m\ge0$ of the regular parity for a nonempty $J$, write
\[
 \mathcal V_J(m)=\sum_{\substack{0\le r\le m\\r\equiv m\ (2)}}
                   \binom mr2^{r-m}M_{J,r}
          =m![t^m]\bigl(\cosh(t/2)G_J(t)\bigr),
 \qquad \mathcal V_\varnothing(m)=0.
\]
For a label $i\in I$, with $J=I\setminus\{i\}$, define
\begin{equation}\label{mtp:W}
 \mathcal W_i(k)=
 \begin{cases}
  \frac12\mathcal S_{J,p_i-k},&k<p_i,\\
  \mathcal V_J(k-p_i)-\frac12\delta_{k,p_i}Z_J,&k\ge p_i.
 \end{cases}
\end{equation}

\begin{theorem}[Triangular mixed-tail recurrence]
\label{mtp:recurrence}
For every $k\ge1$ with $k\equiv d+1\pmod2$,
\begin{equation}\label{mtp:recurrence-formula}
 \boxed{C_{I,k}=\sum_{i\in I}\mathcal W_i(k)
    -\sum_{\substack{S\subsetneq I\\|S|\ge2\text{ even}}}
                     B_{|S|}C_{I\setminus S,k-P_S}.}
\end{equation}
An index $k-P_S\le0$ means the explicit anchor
$\mathcal Z_{I\setminus S,k-P_S}$. Together with
\eqref{mtp:moment-transform}, this is a terminating exact algorithm
for every regular moment, because every right-side tail product has
strictly smaller cardinality.
\end{theorem}

\begin{proof}
Multiply \eqref{mtp:polarized-formula} by $n^k$, sum through a finite
$N$, and then take constant coefficients. When $m=k-p_i<0$, the
symmetric sum is absolutely convergent and equals
$\mathcal S_{J,-m}$. When $m\ge0$, changing the index in one of the
two finite sums gives
\begin{align*}
 \sum_{n=1}^N n^m\bigl(F_J(n-1)+F_J(n)\bigr)
 ={}&\sum_{n=0}^{N-1}\bigl((n+1)^m+n^m\bigr)F_J(n)\\
    &-\delta_{m,0}Z_J+N^mF_J(N).
\end{align*}
The centered polynomial $(x+1/2)^m+(x-1/2)^m$ has the regular parity
for $J$, so the constant coefficient of the first term is
$2\mathcal V_J(m)$. This remains correct for $J=\varnothing$, where
that coefficient is zero by finite power summation.
The difference terms yield $C_{I\setminus S,k-P_S}$ minus their
boundary coefficients. The complete leftover boundary is precisely
\eqref{mtp:boundary-coefficient} and vanishes. What remains is
\eqref{mtp:recurrence-formula}. In particular the exceptional endpoint
term $-Z_J/2$ at $k=p_i$ has not been discarded.
\end{proof}

Let $I^q$ denote the anchored ordinary $q$-fold primitive
\[
 (I^qf)(t)=\frac1{(q-1)!}\int_0^t(t-u)^{q-1}f(u)\,du.
\]
Here the symbol $I^q$ is an integration operator, distinguished from
the label set $I$ by its superscript. Define the finite polynomial
\begin{align}
 R_I(t)={}&\frac12\sum_{i\in I}
       \sum_{\substack{1\le k<p_i\\k\equiv d+1\ (2)}}
          \mathcal S_{I\setminus\{i\},p_i-k}\frac{t^k}{k!}
       -\frac{\mathbf1_{d\ \mathrm{odd}}}{2}
          \sum_{i\in I}Z_{I\setminus\{i\}}\frac{t^{p_i}}{p_i!}
       \notag\\
 &-\sum_{\substack{S\subsetneq I\\|S|\ge2\text{ even}}}B_{|S|}
       \sum_{\substack{1\le k\le P_S\\k\equiv d+1\ (2)}}
          \mathcal Z_{I\setminus S,k-P_S}\frac{t^k}{k!}.
       \label{mtp:R}
\end{align}

\begin{theorem}[Anchored generators, finite arithmetic span, and polylogarithms]
\label{mtp:generator}
The series in \eqref{mtp:GH} converge for $|t|<2\pi$, and are given by
\begin{equation}\label{mtp:generator-formula}
 \boxed{\begin{aligned}
 H_I(t)={}&R_I(t)
       +\sum_{i\in I}I^{p_i}\bigl(\cosh(t/2)G_{I\setminus\{i\}}(t)\bigr)\\
       &-\sum_{\substack{S\subsetneq I\\|S|\ge2\text{ even}}}
                     B_{|S|}I^{P_S}H_{I\setminus S}(t),
       \qquad G_I(t)=\frac{H_I(t)}{2\sinh(t/2)}.
 \end{aligned}}
\end{equation}
For fixed $I$, every $M_{I,m}$ lies in the rational span of $1$ and
convergent positive-index MZVs of weight at most
\begin{equation}\label{mtp:weight-bound}
                         P_I-1-\mathbf1_{d\ \mathrm{odd}}.
\end{equation}
The spanning list is finite and independent of $m$.
Moreover $H_I$ has a finite representation by polynomials in $t$,
MZV constants, and one-variable multiple polylogarithms at $e^{-t}$
with functional depth at most $d-1$. This is a representation bound,
not a claim about minimal depth or arithmetic independence.
\end{theorem}

\begin{proof}
Summing \eqref{mtp:recurrence-formula} with $t^k/k!$ gives
\eqref{mtp:generator-formula} formally. The ranges $k<p_i$, the
endpoint $k=p_i$, and the shifted indices $k-P_S\le0$ give exactly
\eqref{mtp:R}. The remaining positive shifted indices give the displayed
anchored primitives.

Induction on $d$ proves analyticity. All numerators are holomorphic in
$|t|<2\pi$, have parity $d+1$, and vanish to order at least one at
even $d$, or at least two at odd $d$. The denominator has a simple
zero at zero and no other zero in this disk. Its singularity is
therefore removable, and the formal moment series converges there.

Every non-endpoint anchor in $R_I$ has weight $P_I-k$, and its least
possible positive $k$ is one at even $d$ and two at odd $d$. An
endpoint product has weight $P_I-p_i\le P_I-2$ and only occurs at
odd $d$; expand it into MZVs by ordered set partitions. All other
arithmetic constants already belong to generators on proper subsets.
The remaining Taylor coefficients, integrations, and division by
$2\sinh(t/2)$ have rational coefficients. This proves the sharper
parity-sensitive bound \eqref{mtp:weight-bound}, as well as finiteness.
The list can be redundant.

For the functional statement, start at real $t>0$, put $z=e^{-t}$,
and use
\[
 \cosh(t/2)G_J(t)=\left(\frac12+\frac z{1-z}\right)H_J(t).
\]
For $L_{\boldsymbol a}(t)=\operatorname{Li}_{a_1,\ldots,a_\ell}(e^{-t})$
use the convention with $e^{-tn_1}$ on the largest summation index.
Maintain the induction invariant $a_1\ge2$ for nonempty words.
Ordinary integration raises that first index. An explicit unanchored
$q$-fold primitive of $t^m zL_{\boldsymbol a}(t)/(1-z)$ is
\begin{equation}\label{mtp:polylog-primitive}
 (-1)^q\sum_{j=0}^m\binom{q+j-1}{j}\frac{m!}{(m-j)!}
                         t^{m-j}L_{q+j,\boldsymbol a}(t).
\end{equation}
It follows by differentiation from
$L_{1,\boldsymbol a}'=-zL_{\boldsymbol a}/(1-z)$ and
$L_{r+1,\boldsymbol a}'=-L_{r,\boldsymbol a}$. Here $q=p_i\ge2$,
so the new first index is at least two, and depth increases by at
most one.

Subtract the Taylor polynomial of degree $q-1$ at zero from the
\emph{assembled} primitive. Endpoint values with first index at least
two are convergent MZVs. The only possible logarithm comes from a
single first index one. For admissible $\boldsymbol a$,
\begin{align*}
 \lim_{t\to0^+}\bigl(L_{1,\boldsymbol a}(t)
                  -\zeta(\boldsymbol a)L_1(t)\bigr)
 ={}&-\sum_{i=1}^{\ell}\zeta(a_1,\ldots,a_i,1,a_{i+1},\ldots,a_\ell)\\
    &-\sum_{i=1}^{\ell}\zeta(a_1,\ldots,a_i+1,\ldots,a_\ell).
\end{align*}
This follows by subtracting $\zeta(\boldsymbol a)$ inside the outer
series and ordering its missing tail; the difference is absolutely
summable. For the empty word the difference is zero. The coefficient
of $L_1(t)$ in the full endpoint expression is a multiple of $H_J(0)$,
which is zero. Every integration constant is consequently an explicit
convergent MZV. The singleton generator is a polynomial and has depth
zero; induction gives the bound $d-1$. The combined formula is then
continued by its proved holomorphic generator to the stated disk.
\end{proof}

\paragraph{What has been resolved.}
The theorem handles every product of unequal positive even-order tails
and every regular centered moment. It also sharpens the unrestricted
weight bound by one at odd numbers of factors. It neither changes the
summation convention nor assumes a conjecture about MZV independence.
It does not prove the separate all-unequal-order \emph{bilinear optimal
prefix} conjecture: that problem asks for the nonvanishing of different
specified coefficient matrices.
